\documentclass[10pt,a4paper]{article}

% Compile with: latexmk -pdf -interaction=nonstopmode -halt-on-error abstract_resilient_societies_2027.tex
\usepackage[T1]{fontenc}
\usepackage[utf8]{inputenc}
\usepackage{newtxtext}
\usepackage{microtype}
\usepackage[a4paper,top=24mm,bottom=23mm,left=29mm,right=29mm]{geometry}
\usepackage{xcolor}
\usepackage{parskip}
\usepackage{hyperref}
\usepackage{fancyhdr}

\definecolor{accent}{HTML}{24566A}
\definecolor{ink}{HTML}{1E2931}
\definecolor{muted}{HTML}{65717B}
\hypersetup{hidelinks,pdftitle={Sea Level Rise Weakens the Coastal Protection Value of Wetlands},pdfauthor={Yanxu Long}}
\pagestyle{empty}
\setlength{\parindent}{0pt}
\setlength{\parskip}{0.78em}
\linespread{1.08}
\color{ink}

\begin{document}

{\sffamily\color{accent}\fontsize{8.5}{11}\selectfont\bfseries\MakeUppercase{Resilient Societies}\hspace{0.7em}\textcolor{muted}{\textbar}\hspace{0.7em}\MakeUppercase{Workshop Abstract}}\\[7pt]
{\fontsize{17.5}{20.5}\selectfont\bfseries Sea Level Rise Weakens the Coastal\\[1pt]Protection Value of Wetlands}\par
\vspace{3pt}
{\sffamily\fontsize{9}{12}\selectfont\color{muted}Yanxu Long \hspace{0.5em}\textbar\hspace{0.5em} London School of Economics \hspace{0.5em}\textbar\hspace{0.5em} 20--21 May 2027}\par
\vspace{3pt}
{\color{accent}\rule{\linewidth}{0.75pt}}
\vspace{0.5pt}

\begingroup
\fontsize{10.2}{13.3}\selectfont
Coastal wetlands serve as natural infrastructure, attenuating wave energy and reducing storm-related property damage. The economic value of this protection is well documented (Narayan et al., 2017; Sun and Carson, 2020), yet how climate change affects this value remains poorly understood. Rising sea levels may not only intensify coastal hazards but also weaken the natural defenses that protect communities from their impacts. We quantify an important but understudied channel of climate damages: sea-level rise erodes the coastal-protection value of wetlands, potentially amplifying the economic consequences of storms.

We study counties along the U.S. Atlantic and Gulf coasts, combining historical records of storm-related property damage with measures of coastal wetland coverage, pre-storm sea levels, and local characteristics. We extend the empirical framework of Sun and Carson (2020) to examine how sea-level rise affects the coastal-protection value of wetlands. We use two complementary approaches. First, we estimate the protective contribution of wetlands across years with different pre-storm sea levels. Second, we incorporate an interaction between wetland coverage and sea level directly into the storm-damage model to test whether higher seas weaken the association between wetlands and storm losses. Together, these approaches assess whether the benefits people derive from coastal natural capital depend on changing environmental conditions.

Coastal wetlands provide substantial economic protection, with an estimated \$5.7 billion in annual avoided storm-related property damage along the U.S. Atlantic and Gulf coasts (Sun and Carson, 2020). We extend this established valuation by examining how the protective value of wetlands changes with sea level. Across our two empirical approaches, a one-centimeter increase in sea level is associated with a 3.5--10\% reduction in wetland protection, equivalent to approximately \$200--600 million in lost annual protection value. These findings reveal an additional economic cost of sea-level rise: beyond increasing direct flood risks, rising seas can diminish the protective benefits that coastal communities derive from natural capital.

Several ecological mechanisms may explain this relationship. Rising seas can increase wetland inundation, accelerate marsh conversion into open water, and reduce the effectiveness of vegetation and surface roughness in attenuating waves and storm surge. Consequently, even where wetlands persist, greater water depths may weaken their capacity to protect coastal communities from extreme events.

Our findings contribute to the economics of ecosystem services, natural capital valuation, and climate adaptation by showing that the benefits people receive from wetlands are sensitive to sea-level change. The erosion of natural protection may also have unequal consequences: communities with fewer resources to replace declining natural defenses may face greater adaptation burdens. More broadly, climate change may simultaneously intensify environmental shocks and weaken the natural capital that buffers their consequences. Accounting for these changing ecosystem-service values is therefore important for evaluating coastal resilience and designing adaptation strategies.
\endgroup

\end{document}
